\chapter{\ifproject%
\ifenglish Project Structure and Methodology\else โครงสร้างและขั้นตอนการทำงาน\fi
\else%
\ifenglish Project Structure\else โครงสร้างของโครงงาน\fi
\fi
}

ในบทนี้จะกล่าวถึงภาพรวมของระบบ หลักการ และการออกแบบแต่ละขั้นตอนของ Context Construction Pipeline
ตั้งแต่การรับ Alert และ Raw Logs จาก SIEM จนถึงการสร้างรายงานการสืบสวน
โดยใช้ชุดข้อมูล Account Discovery: Local Account (T1087.001) จาก Splunk Attack Data เป็นตัวอย่างประกอบ

\section{ภาพรวมของระบบ}

ภาพรวมของระบบแสดงดังรูปที่~\ref{fig:overview} แบ่งเป็น 4 ส่วน ได้แก่
\textbf{Input} (Alert และ Raw Logs),
\textbf{Context Construction Pipeline} (Normalize, Event Selection, Representation),
\textbf{Analysis} (RAG และ LLM Analysis) และ \textbf{Output} (Investigation Report)

\begin{figure}[ht]
\centering
\resizebox{\textwidth}{!}{%
\begin{tikzpicture}[
  node distance=6mm and 7mm,
  box/.style={draw, rounded corners=2pt, fill=white, align=center,
              minimum width=23mm, minimum height=13mm, font=\small},
  arr/.style={-{Latex[length=2.5mm]}, very thick, blue!50!black},
  grp/.style={draw, dashed, rounded corners=8pt, gray, inner sep=3mm}]
  \node[box] (alert) {Alert/EDR\\[-1pt]{\scriptsize ตรวจพบความผิดปกติ}};
  \node[box, below=of alert] (raw) {Raw Log\\[-1pt]{\scriptsize นำ Log จาก SIEM}};
  \node[box, right=12mm of raw] (norm) {Normalize\\[-1pt]{\scriptsize CIM standard}};
  \node[box, right=of norm] (sel) {Event Selection\\[-1pt]{\scriptsize คัด Event สำคัญ}};
  \node[box, right=of sel] (rep) {Representation\\[-1pt]{\scriptsize JSON / TOON}};
  \node[box, right=12mm of rep] (llm) {LLM Analysis\\[-1pt]{\scriptsize ระบุ Technique}};
  \node[box, above=of llm] (rag) {RAG\\[-1pt]{\scriptsize ดึง Top-$k$ MITRE}};
  \node[box, right=12mm of llm] (out) {Output\\[-1pt]{\scriptsize Report}};
  \draw[arr] (alert) -- (raw);
  \draw[arr] (raw) -- (norm);
  \draw[arr] (norm) -- (sel);
  \draw[arr] (sel) -- (rep);
  \draw[arr] (rep) -- (llm);
  \draw[arr] (rep.east) ++(4mm,0) |- (rag.west);
  \draw[arr] (rag) -- (llm);
  \draw[arr] (llm) -- (out);
  \begin{scope}[on background layer]
    \node[grp, fit=(alert)(raw), label={[gray]above:Input}] {};
    \node[grp, fit=(norm)(sel)(rep), label={[gray]above:Context Construction Pipeline}] {};
    \node[grp, fit=(rag)(llm), label={[gray]above:Analysis}] {};
    \node[grp, fit=(out), label={[gray]above:Output}] {};
  \end{scope}
\end{tikzpicture}}
\caption{ภาพรวมของระบบและ Context Construction Pipeline}
\label{fig:overview}
\end{figure}

หลักการสำคัญของระบบคือ \textbf{ระบบไม่ได้ทำหน้าที่ตรวจจับภัยคุกคามแทน SIEM/EDR}
Alert ที่ระบบเดิมตรวจพบเป็นเพียงจุดเริ่มต้นการสืบสวน (Input ที่รับเข้ามา)
หน้าที่ของ Pipeline คือขยายข้อมูลรอบจุดเริ่มต้นนั้นให้เป็น Investigation Context
ที่ LLM สามารถวิเคราะห์ได้อย่างมีหลักฐาน
การจัดระเบียบข้อมูลก่อนส่งให้ LLM ช่วยลด Noise และทำให้ LLM โฟกัสข้อมูลที่สำคัญได้ดีกว่าการส่ง Raw Logs โดยตรง
Input และ Output ของแต่ละขั้นตอนสรุปได้ดังตารางที่~\ref{tab:pipeline-io}

\begin{table}[ht]
  \centering
  \caption{Input และ Output ของแต่ละขั้นตอนใน Pipeline}
  \label{tab:pipeline-io}
  \begin{tabular}{|l|l|l|}
    \hline
    \textbf{ขั้นตอน} & \textbf{Input} & \textbf{Output} \\ \hline
    Raw Log        & Alert และคำค้นไปยัง SIEM            & Raw Logs (XML/JSON) \\ \hline
    Normalize      & Windows/Linux Raw Logs               & JSON ตามโครงสร้าง Field ร่วม (CIM) \\ \hline
    Event Selection & JSON Logs + Alert                   & Selected JSON Logs \\ \hline
    Representation & Selected JSON Logs                   & Representation Log \\ \hline
    RAG            & Query จาก Representation Log         & Top-$k$ Techniques (Context) \\ \hline
    LLM Analysis   & Context + Representation Log         & Attack Technique พร้อมหลักฐาน \\ \hline
    Output         & ผลการวิเคราะห์ของ LLM                 & Investigation Report \\ \hline
  \end{tabular}
\end{table}

\section{ขั้นตอนที่ 1: การรับ Alert และ Raw Logs จาก SIEM}

Logs จากชุดข้อมูลถูก Replay เข้า Splunk Enterprise เพื่อจำลองสภาพแวดล้อม SIEM จริง
เมื่อมี Alert เกิดขึ้น (หรือในการทดลองใช้ Case Configuration จำลอง Alert)
Pipeline จะดึง Logs ที่เกี่ยวข้องจาก Splunk ผ่าน REST API และบันทึกเป็นไฟล์ JSON Lines
สำหรับชุดข้อมูลตัวอย่าง T1087.001 ซึ่งเป็น Windows Sysmon Logs ใน Splunk
มีจำนวน 11,090 events ในช่วงเวลาประมาณ 15 นาที (ขนาดไฟล์ XML ประมาณ 16.826 MB)

\section{ขั้นตอนที่ 2: Normalize}

Raw Logs จากแต่ละแหล่งมีรูปแบบต่างกัน เช่น Sysmon ใน Splunk อยู่ในรูป \texttt{XmlWinEventLog}
และ Auditd เป็นข้อความหลาย Record ต่อเหตุการณ์
ขั้นตอน Normalize จะ Parse Logs แล้วแปลงเป็น JSON ที่มีโครงสร้าง Field ร่วม (Common Field Structure)
ตามแนวทางของ Splunk CIM~\cite{splunk-cim} เช่น \texttt{EventID}, \texttt{SystemTime}, \texttt{RecordNumber},
\texttt{Channel}, \texttt{Image}, \texttt{CommandLine}, \texttt{ProcessGuid}, \texttt{ParentProcessGuid} และ \texttt{User}
ดังตัวอย่างในรูปที่~\ref{fig:normalize}
ผลลัพธ์ทำให้ขั้นตอนถัดไปทำงานกับ Logs จากทั้ง Windows และ Linux ได้ด้วยตรรกะเดียวกัน

\begin{figure}[ht]
  \centering
  \includegraphics[width=\textwidth]{normalize_example.png}
  \caption{ตัวอย่าง Raw Logs (XML) และผลลัพธ์หลัง Normalize เป็น JSON}
  \label{fig:normalize}
\end{figure}

\section{ขั้นตอนที่ 3: Event Selection}

Event Selection ตอบคำถามว่า \emph{``จะส่งเหตุการณ์ใดให้ LLM''}
ซึ่งเป็นขั้นตอนที่ยากที่สุดของ Pipeline เพราะต้องคัดเลือกให้ได้หลักฐานครบ แต่ไม่มี Noise มากเกินไป
โครงงานเสนอวิธี Event Selection 5 วิธีเพื่อนำมาเปรียบเทียบ (รูปที่~\ref{fig:es-methods} และตารางที่~\ref{tab:es-methods})

\begin{figure}[ht]
  \centering
  \includegraphics[width=\textwidth]{event_selection_methods.png}
  \caption{แนวคิดของวิธี Event Selection แบบที่ 1--4 (แถบสีเขียวคือ Event ที่ถูกเลือก)}
  \label{fig:es-methods}
\end{figure}

\begin{enumerate}
  \item \textbf{Alert Only} -- ใช้เฉพาะข้อมูลจาก Alert
  \item \textbf{Single Event} -- ดึงเฉพาะ Event เดียวที่ Alert ชี้
  \item \textbf{Time Window} -- ดึง Events ในช่วงเวลา $\pm T$ รอบ Alert
  \item \textbf{Multiple Related Events} -- ดึง Events ที่เชื่อมโยงกันจริงผ่านความสัมพันธ์ของ Process
        (ProcessGuid Lineage) ได้แก่ Ancestor, Descendant และ Events อื่นที่ ProcessGuid เดียวกัน
  \item \textbf{Top-$k$ Relevant Events} -- ให้คะแนนความเกี่ยวข้องกับทุก Event ในขอบเขตที่กำหนด
        แล้วเลือกเฉพาะ Events ที่มีคะแนนสูงสุดภายใต้งบประมาณ Token
\end{enumerate}

\begin{table}[ht]
  \centering
  \caption{ข้อดีและข้อเสียของวิธี Event Selection แต่ละแบบ}
  \label{tab:es-methods}
  \begin{tabular}{|L{0.16\textwidth}|L{0.36\textwidth}|L{0.36\textwidth}|}
    \hline
    \textbf{วิธี} & \textbf{ข้อดี} & \textbf{ข้อเสีย} \\ \hline
    Alert Only &
      ประมวลผลเร็วที่สุด ไม่มี Overhead จากการ Query และใช้ได้แม้ไม่มี Raw Logs &
      ขาดหลักฐานเชิงลึก ไม่สามารถทำ Evidence Grounding ได้ LLM ได้แต่เชื่อตาม Alert \\ \hline
    Single Event &
      ได้รายละเอียดเต็มของ Event (CommandLine, GUID) และแยกผลของ ``รายละเอียด 1 Event'' ได้ชัด &
      ไม่เห็นว่า Event มาจากไหนหรือต่อไปทำอะไร พลาด Attack Chain ที่กระจายหลาย Event \\ \hline
    Time Window &
      Implement ง่าย (Range Query) ใช้ได้กับ Logs ทุกประเภท และจับ Events ที่เกิดพร้อมกันแต่ไม่เชื่อมกันทางโครงสร้างได้ &
      Noise สูงถ้า $T$ กว้าง พลาด Event ถ้า $T$ แคบ และไม่มีค่า $T$ ที่ถูกต้องสากล \\ \hline
    Multiple Related &
      สร้าง Attack Chain ได้ และตัด Events ที่เกิดพร้อมกันโดยบังเอิญออก &
      พึ่งพา Relationship Key สูง อาจพลาดพฤติกรรมที่ไม่เชื่อมทาง Process และเสี่ยง Over-traversal หากไม่จำกัดความลึก \\ \hline
    Top-$k$ Relevant &
      ควบคุมขนาด Context ได้ (Predictable Context Size) ยืดหยุ่น ปรับ $k$ ตามข้อจำกัด Token ได้ &
      ต้องออกแบบ Scoring Model ให้เป็น Generic และอ่อนไหวต่อค่า $k$ \\ \hline
  \end{tabular}
\end{table}

จากตัวอย่าง T1087.001 เมื่อกำหนดให้ Alert ชี้ที่ Process \texttt{net.exe} (คำสั่ง \texttt{net users})
และใช้วิธี Multiple Related เพื่อดึง Process ที่เกี่ยวข้อง (ผู้ใช้ที่รัน, Parent และ Child Process)
พร้อม Events ที่ Process เหล่านั้นเป็นเจ้าของ จำนวน Events ลดลงจาก 11,090 เหลือ 177 events
วิธี Top-$k$ Relevant Events ซึ่งเป็นวิธีที่โครงงานนำเสนอ ต่อยอดจากแนวคิดนี้โดยแบ่งเป็น 4 ขั้นตอนย่อย ดังนี้

\subsection{Graph: สร้างกราฟความสัมพันธ์ของ Process}
สร้างกราฟความสัมพันธ์ระหว่าง Process และ Event โดยมีโหนดเป็น Process (ระบุด้วย ProcessGuid)
และเส้นเชื่อม 3 ประเภท ได้แก่
(1) Parent $\rightarrow$ Child จาก Sysmon EID 1 (Process Creation)
(2) Process Access จาก EID 10 และ
(3) File $\rightarrow$ Exec จาก EID 11 (ไฟล์ที่ถูกสร้างแล้วถูกรัน)

\subsection{Traverse: ท่องกราฟจากจุดที่ Alert ชี้}
เริ่มจาก Process ที่ Alert ชี้ (Anchor) แล้วขยายไปยัง Parent และ Child Process
โดยเก็บเฉพาะบรรพบุรุษและลูกหลานโดยตรงที่อยู่ในระยะที่กำหนด
เพื่อให้ได้ Candidate Events ที่เกี่ยวข้องกับ Alert โดยไม่ลามไปทั้งระบบ ตามกฎดังนี้
\begin{enumerate}
  \item เดินตามความสัมพันธ์ Parent $\rightarrow$ Child
  \item เดินทางเดียว คือขึ้นตรงหรือลงตรง ไม่เลี้ยวเข้ากิ่งข้างเคียง
  \item หยุดเมื่อครบจำนวน Hop สูงสุด หรือเจอ Boundary
        (Process ที่รู้ว่ามีอยู่ แต่ไม่มี EID 1 หรือบันทึกการเกิดในไฟล์ Log)
  \item ทุก Process ที่เดินถึง จะเก็บทุก Event ที่ Process นั้นเป็นเจ้าของ
\end{enumerate}
รูปที่~\ref{fig:traverse} แสดงส่วนหนึ่งของกราฟที่ได้จาก Sysmon Log ของ T1087.001
โดยเส้นทึบสีดำคือความสัมพันธ์ Parent--Child (EID 1) และเส้นประสีแดงคือ Process Access (EID 10)

\begin{figure}[ht]
  \centering
  \includegraphics[width=0.85\textwidth]{traverse_graph.png}
  \caption{ส่วนหนึ่งของกราฟ Process จาก Sysmon Log ของ T1087.001 โดยมี \texttt{net.exe} เป็นจุดเริ่ม (Anchor)}
  \label{fig:traverse}
\end{figure}

\subsection{Score: ให้คะแนนความเกี่ยวข้อง}
ประเมินและจัดลำดับความเกี่ยวข้องของแต่ละ Candidate Event $e$ ด้วยผลคูณของ 3 ปัจจัย ดังสมการ~\eqref{eq:score}

\begin{equation}\label{eq:score}
\mathrm{Score}(e) = \mathrm{Prox}(e) \times \mathrm{Type}(e) \times \mathrm{Time}(e)
\end{equation}
\begin{align}
\mathrm{Prox}(e) &= \gamma^{\,hop(e)} \\
\mathrm{Type}(e) &= w[\mathrm{EventID}] \\
\mathrm{Time}(e) &= e^{-|\Delta t|/\tau}
\end{align}

\noindent
โดยที่ $hop(e)$ คือจำนวนขั้นบนกราฟจากจุดที่ Alert ชี้ถึง Process เจ้าของ Event,
$\gamma$ คือสัดส่วนคะแนนที่เหลือเมื่อห่างออกไป 1 Hop,
$w[\mathrm{EventID}]$ คือน้ำหนักของ Event แต่ละชนิด (0--1) ที่บอกว่า Event ชนิดนั้นช่วยการวิเคราะห์ได้มากเพียงใด,
$\Delta t$ คือเวลาที่ห่างจาก Alert (วินาที) และ $\tau$ คือช่วงเวลาที่คะแนนลดลงเหลือประมาณ $1/e$ (ราว $1/3$)
ค่าพารามิเตอร์ที่ใช้ในปัจจุบันแสดงในตารางที่~\ref{tab:score-param}
ซึ่งยังเป็นค่าตั้งต้นที่ต้องปรับจูนในการทดลองขั้นต่อไป
ตัวอย่างการคำนวณคะแนนแสดงในตารางที่~\ref{tab:score-example}

\begin{table}[ht]
  \centering
  \caption{ค่าพารามิเตอร์ตั้งต้นของการให้คะแนน}
  \label{tab:score-param}
  \begin{tabular}{|c|c|l|}
    \hline
    \textbf{พารามิเตอร์} & \textbf{ค่า} & \textbf{ความหมาย} \\ \hline
    $\gamma$ & 0.6 & Hop 0 = 1, Hop 1 = 0.6, Hop 2 = 0.36 \\ \hline
    $w[\text{EID 1}]$ & 1.0 & Process Creation มีข้อมูลคำสั่งครบ \\ \hline
    $w[\text{EID 10}]$ & 0.2 & Process Access ให้ข้อมูลประกอบ \\ \hline
    $\tau$ & 300 วินาที & ห่าง 300 วินาที คะแนนเวลาเหลือ 0.37 \\ \hline
  \end{tabular}
\end{table}

\begin{table}[ht]
  \centering
  \caption{ตัวอย่างการคำนวณคะแนนของ Candidate Events จาก T1087.001}
  \label{tab:score-example}
  \begin{tabular}{|l|c|c|c|c|c|c|}
    \hline
    \textbf{Event} & \textbf{Hop} & \textbf{$\Delta t$} & \textbf{Prox} & \textbf{Type} & \textbf{Time} & \textbf{Score} \\ \hline
    net1.exe (EID 1)         & 1 & 0.01 s & 0.6 & 1.0 & 1.00 & 0.60 \\ \hline
    cmd.exe (EID 1)          & 1 & 6.6 s  & 0.6 & 1.0 & 0.98 & 0.59 \\ \hline
    net $\rightarrow$ net1 (EID 10) & 0 & 0 s & 1.0 & 0.2 & 1.00 & 0.20 \\ \hline
    cmd $\rightarrow$ WMIC (EID 10) & 1 & 11 s & 0.6 & 0.2 & 0.97 & 0.12 \\ \hline
    cmd $\rightarrow$ cmd (EID 10)  & 1 & 74 s & 0.6 & 0.2 & 0.78 & 0.09 \\ \hline
  \end{tabular}
\end{table}

\subsection{Budget: คัดเลือกภายใต้งบประมาณ Token}
เพื่อจำกัดจำนวน Token ให้พอดีกับ Context Window ของโมเดล และให้ LLM เห็นเฉพาะ Events ที่สำคัญที่สุดก่อน
ขั้นตอน Budget ทำงาน 3 ข้อ ดังนี้
\begin{enumerate}
  \item ใส่ Event ที่ Alert ชี้ (Anchor) เป็นอันดับแรก
  \item เรียง Events ที่เหลือตามคะแนนจากมากไปน้อย
  \item ใส่ทีละ Event พร้อมนับจำนวน Token สะสม หากเกินโควตา $B$ ให้หยุด
\end{enumerate}
ตัวอย่างเช่น หากกำหนดโควตา $B = 400$ Tokens และแต่ละ Event ใช้ประมาณ 100--130 Tokens
จะมีเพียงประมาณ 3 Events แรกเท่านั้นที่ถูกส่งต่อไปยังขั้นตอนถัดไป

\section{ขั้นตอนที่ 4: Representation}

Representation ตอบคำถามว่า \emph{``จะนำเสนอ Logs ที่คัดเลือกแล้วให้ LLM ในรูปแบบใดจึงเหมาะสมที่สุด''}
เนื่องจากข้อมูลชุดเดียวกันเขียนต่างรูปแบบกันจะได้จำนวน Token และความเข้าใจของ LLM ที่ต่างกัน
โครงงานเปรียบเทียบ 5 รูปแบบ (รูปที่~\ref{fig:representation} และตารางที่~\ref{tab:representation})
โดยพิจารณาตามลำดับความสำคัญ คือ (1) Accuracy ในการระบุเทคนิค และ (2) Token Efficiency

\begin{figure}[ht]
  \centering
  \includegraphics[width=\textwidth]{representation_formats.png}
  \caption{ตัวอย่าง Event เดียวกันในรูปแบบ JSON, XML, TOON, Key=Value และ Timeline}
  \label{fig:representation}
\end{figure}

\begin{table}[ht]
  \centering
  \caption{จำนวน Token โดยประมาณของ Event ตัวอย่าง (Sysmon EID 1) ในแต่ละรูปแบบ}
  \label{tab:representation}
  \begin{tabular}{|l|c|l|}
    \hline
    \textbf{รูปแบบ} & \textbf{Tokens (ประมาณ)} & \textbf{ลักษณะ} \\ \hline
    Timeline   & 35  & บรรทัดเดียว กระชับที่สุด แต่ข้อมูลบางส่วนหายไป \\ \hline
    TOON       & 119 & ตัด \{\} และ "" ออกจาก JSON \\ \hline
    Key=Value  & 120 & รูปแบบที่ Splunk แสดงผล \\ \hline
    JSON       & 138 & Output ของขั้นตอน Normalize (ใช้ในการทดลองปัจจุบัน) \\ \hline
    XML        & 195 & รูปแบบต้นฉบับของ Sysmon \\ \hline
  \end{tabular}
\end{table}

ในการทดลองปัจจุบันโครงงานเลือกใช้รูปแบบ JSON เนื่องจากได้มาโดยตรงจากขั้นตอนก่อนหน้า
ส่วนการเปรียบเทียบ Accuracy ระหว่างรูปแบบต่าง ๆ จะดำเนินการในการทดลองขั้นต่อไป

\section{ขั้นตอนที่ 5: Retrieval-Augmented Generation (RAG)}

RAG ทำหน้าที่เพิ่มบริบทให้ LLM โดยอ้างอิงองค์ความรู้จาก MITRE ATT\&CK
ฐานความรู้สร้างจากเอกสาร MITRE ATT\&CK Enterprise v19.2 (STIX ทางการ)
แปลงเป็น JSON และแบ่ง Chunk ตาม Technique และ Sub-technique
(1 Chunk ต่อ 1 Technique หรือ 1 Sub-technique) รวมทั้งหมด 697 Chunks (222 Techniques + 475 Sub-techniques)
ดังรูปที่~\ref{fig:rag-kb}
ข้อความของแต่ละ Chunk ประกอบด้วยชื่อ Tactics และคำอธิบาย (Description) โดยไม่ Embed รหัสเทคนิค
เหตุผลที่แบ่งถึงระดับ Sub-technique คือคำตอบที่ต้องการอยู่ในระดับ Sub-technique
และ Sub-technique ภายใต้ Technique เดียวกันมักต่างกันเพียงในคำอธิบาย
เช่น T1087.001 (\texttt{net user}) กับ T1087.002 (\texttt{net user /domain})
แต่ละ Chunk ถูกแปลงเป็นเวกเตอร์ด้วย \texttt{qwen3-embedding:0.6b} ผ่าน Ollama
ความยาวมัธยฐานของ Chunk ประมาณ 1,100 ตัวอักษร (ประมาณ 280 Tokens)

\begin{figure}[ht]
  \centering
  \includegraphics[width=\textwidth]{rag_kb_chunks.png}
  \caption{ตัวอย่างฐานความรู้ MITRE ATT\&CK v19.2 ที่แบ่ง Chunk ตาม Technique และ Sub-technique}
  \label{fig:rag-kb}
\end{figure}

Query ที่ใช้ค้นหาในฐานความรู้มีได้ 2 แบบ คือ
(1) ใช้ Representation Log (JSON) โดยตรง และ
(2) ให้ LLM แปลง Representation Log เป็นคำอธิบายพฤติกรรมภาษาอังกฤษ (Log Summary) ก่อน
โดยห้ามใช้ชื่อหรือรหัสเทคนิค ผลลัพธ์ของ RAG คือรายการ Top-$k$ Techniques ที่ใช้เป็น Context ให้ LLM

\section{ขั้นตอนที่ 6: LLM Analysis และ Output}

LLM ที่ใช้คือ Qwen3.5:9B ทำงานภายในเครื่องผ่าน Ollama ร่วมกับ Hermes เป็น Harness
LLM รับ Context จาก RAG และ Representation Log แล้วระบุเทคนิคการโจมตี
Prompt ถูกออกแบบให้ LLM ทำงานเป็นลำดับ ได้แก่
อ่าน Event ที่ Alert ชี้ก่อน $\rightarrow$ พิจารณาว่าคำสั่งทำอะไร $\rightarrow$ ระบุ Technique
$\rightarrow$ ระบุ Sub-technique ที่ตรงกับรูปแบบคำสั่ง $\rightarrow$ ระบุทางเลือกที่ใกล้เคียงและเหตุผลที่ไม่ใช่
และตอบในรูปแบบ JSON ที่มีโครงสร้างกำหนด (เช่น พฤติกรรมที่สังเกตได้ ทางเลือกที่พิจารณา รหัสเทคนิค และความมั่นใจ)
สำหรับกรณีใช้ RAG จะระบุใน Prompt ว่ารายการ Candidate เป็นเพียงคำใบ้ และคำตอบที่ถูกอาจไม่อยู่ในรายการ
การตั้งค่าในการทดลอง ได้แก่ ปิดโหมด Thinking, Temperature = 0 และ Context Window 16,384 Tokens

เพื่อป้องกันการรั่วไหลของเฉลย ระบบจะตรวจสอบว่า Prompt และ Query ไม่มีรหัสเทคนิค
(รูปแบบ \texttt{T\#\#\#\#}) หรือรหัส Case ที่บอกคำตอบ ก่อนส่งให้ LLM
ผลลัพธ์สุดท้ายคือ Investigation Report ที่ระบุเทคนิค ATT\&CK หลักฐาน (Event ที่อ้างอิง) และเหตุผลประกอบ

\section{ต้นแบบ Web Dashboard}

ผลการวิเคราะห์ถูกนำเสนอให้ SOC Analyst ผ่าน Web Dashboard
ซึ่งรองรับการรับข้อมูล 2 แบบ คือ ดึง Alert จาก Splunk อัตโนมัติ และอัปโหลดไฟล์ Logs ด้วยตนเอง (Manual Investigation)
หน้า Inbox (รูปที่~\ref{fig:dashboard}) แสดงรายการการสืบสวนทั้งหมด พร้อมแหล่งที่มา ความรุนแรง เวลา
ชื่อ Alert และ Host สถานะการประมวลผลของ LLM (Pending, Processing, Analyzed, Failed)
ผลลัพธ์ (เทคนิคที่ระบุได้) และผู้รับผิดชอบ
นอกจากนี้ยังมีเมนูสำหรับผู้ดูแลระบบ เช่น การเชื่อมต่อ SIEM การตั้งค่า Pipeline
และการปรับปรุงข้อมูล/ฐานความรู้ ซึ่งกำหนดให้ผู้ดูแลระบบเป็นผู้สั่งปรับปรุงเท่านั้น

\begin{figure}[ht]
  \centering
  \includegraphics[width=0.95\textwidth]{dashboard_prototype.png}
  \caption{ต้นแบบหน้า Inbox ของ Web Dashboard}
  \label{fig:dashboard}
\end{figure}
